\documentclass[11pt]{article}
\usepackage[T1]{fontenc}
\usepackage[utf8]{inputenc}
\usepackage{lmodern,amsmath,amssymb,amsthm,mathtools}
\usepackage[margin=0.86in]{geometry}
\usepackage{microtype,booktabs,longtable,array,enumitem}
\usepackage{xcolor,hyperref,fancyhdr}
\definecolor{ink}{HTML}{153A56}
\hypersetup{colorlinks=true,linkcolor=ink,urlcolor=ink,citecolor=ink,
pdftitle={Rogers--Szego discriminants: removing the cyclotomic guards},
pdfauthor={Research working note}}
\pagestyle{fancy}\fancyhf{}\fancyhead[L]{\small Rogers--Szeg\H{o} discriminants}
\fancyhead[R]{\small 17 October 2026}\fancyfoot[C]{\thepage}
\setlength{\headheight}{14pt}
\setlength{\parskip}{4pt}
\setlength{\parindent}{0pt}
\setlist{nosep,leftmargin=1.5em}
\newtheorem{theorem}{Theorem}
\newtheorem{lemma}[theorem]{Lemma}
\newtheorem{corollary}[theorem]{Corollary}
\theoremstyle{remark}\newtheorem{remark}[theorem]{Remark}
\newcommand{\He}{\operatorname{He}}
\newcommand{\Disc}{\operatorname{Disc}}
\newcommand{\Res}{\operatorname{Res}}
\newcommand{\ord}{\operatorname{ord}}
\newcommand{\Q}{\mathbb{Q}}
\newcommand{\Z}{\mathbb{Z}}
\newcommand{\eps}{\varepsilon}
\newcommand{\F}{\mathbb{F}}
\newcommand{\poch}[3]{(#1;#2)_{#3}}
\title{\color{ink}Rogers--Szeg\H{o} discriminants:\\
removing the cyclotomic guards}
\author{Extension of the existing local-discriminant research}
\date{17 October 2026}
\begin{document}
\maketitle
\begin{abstract}
The previous local ordinary-discriminant theorem required residual
squarefreeness, separation, and simplicity of an auxiliary Hermite
polynomial. We remove the auxiliary condition for every odd root order
and every odd residual degree. We also remove the static conditions for
all residual degrees at every odd prime order. Consequently the exact
prime cyclotomic order of the Rogers--Szeg\H{o} ordinary discriminant is
unconditional, in all degrees. The new arithmetic step proves that every
root of an adjacent parity-Hermite Wronskian satisfies
$v_2(y-1)=1/2$, excluding collisions over odd cyclotomic fields.
The proof extends the earlier theorem; its local expansion and leading
coefficient formulas are retained as established inputs.
This is mathematical progress in the algebraic branch. It supplies no
improved FFT complexity bound. Historical priority beyond the comparisons
recorded here remains unresolved.
\end{abstract}

\section{The result and the exact gain}

Define
\[
H_n(z;q)=\sum_{k=0}^n \genfrac{[}{]}{0pt}{}{n}{k}_q z^k,\qquad
D_n(q)=\Disc_z H_n(z;q)\in\Z[q].
\]
Gaussian binomial coefficients are evaluated as polynomials, including
at roots of unity. All discriminants use the ordinary derivative in
$z$. Our convention is $\Disc(1)=\Disc(\text{monic linear})=1$.

\begin{theorem}[Exact odd-prime cyclotomic orders]\label{thm:main}
Let $p$ be an odd prime and write $n=Ap+b$, where $A\geq0$ and
$0\leq b<p$. Then
\[
\boxed{\ord_{\Phi_p(q)}D_n(q)
       =p\binom{A}{2}+A(b\bmod2).}
\tag{1}
\]
Here the order is exact: one additional factor of $\Phi_p$ does not
divide the remaining quotient. More generally, if $d=p^s$, $s\geq1$,
and $n=Ad+b$ with $0\leq b<p$, then
\[
\boxed{\ord_{\Phi_d(q)}D_n(q)
       =d\binom{A}{2}+A(b\bmod2).}
\tag{2}
\]
For $A\geq1$, the complete leading coefficients in
Section~\ref{sec:leading} are nonzero without extra hypotheses.
\end{theorem}

The prime-power statement deliberately retains $b<p$, rather than the
larger range $b<d$. The theorem includes $A=0$, where the order is zero,
and $n=0,1$ under the conventions above.

The comparison with the previous checkpoint is precise. Version 3 of
\emph{Local ordinary discriminants of Rogers--Szeg\H{o} polynomials at
odd roots of unity}, dated 11 October 2026, assumes the following
nonvanishing conditions:
\[
\begin{array}{ll}
\text{even }b:&\Disc(G)\Res(F,G)\ne0,\\
\text{odd }b:&g\ne0,\quad
\Disc(U)\Res(F,U)\Disc(J_{A+1})\ne0,
\end{array}
\tag{3}
\]
where $F=1+z^d$, $G=H_b(z;\omega)$, $G=(z+1)U$ in the odd case,
and $g=U(-1)$. A later accepted review removed every condition only
for $b=1,3$. The extension here proves:

\begin{center}
\begin{tabular}{>{\raggedright\arraybackslash}p{0.30\linewidth}>{\raggedright\arraybackslash}p{0.59\linewidth}}
\toprule
Condition removed & Proven range\\
\midrule
$g\ne0$ & Every odd $d\geq3$ and odd $1\leq b<d$.\\
$\Disc(J_{A+1})\ne0$ & The same full range, for every $A\geq1$.\\
Residual squarefreeness and separation & Every $d=p^s$ with
odd prime $p$ and $0\leq b<p$.\\
All conditions together & Every odd prime $d=p$, every
$0\leq b<p$, and every $A\geq1$.\\
\bottomrule
\end{tabular}
\end{center}

Thus (1) extends the earlier result to an unconditional theorem
uniform in the prime order, residual degree, and block count.

\section{The overlap parameter for every odd residual}

Fix an odd $d=2h+1\geq3$, a primitive $d$th root $\omega$, and
$b=2r+1<d$. Use the original normalization
\[
W_b=\frac{\poch{\omega}{\omega}{b}}
              {\poch{\omega^2}{\omega^2}{(d+b)/2}},
\quad
\tau_b=\frac{dW_b}{g_b},
\quad
J_{A+1}(x;\tau)
=x\He_A(x)-A\tau\He_{A-1}(x).
\tag{4}
\]
The Hermite polynomials are monic probabilists' Hermite polynomials:
$\He_0=1$, $\He_1=x$, and
$\He_{n+1}=x\He_n-n\He_{n-1}$.

\begin{lemma}[Partial sums and nonreality]\label{lem:partial}
Set
\[
t_k=\frac{\poch{\omega}{\omega^2}{k}}
           {\poch{\omega^2}{\omega^2}{k}},
\qquad S_j=\sum_{k=0}^j t_k.
\]
Then
\[
g_b=\poch{\omega^2}{\omega^2}{r}S_r\ne0,
\qquad
\boxed{\tau_b=\frac{S_h}{S_r}.}
\tag{5}
\]
In the principal embedding $\omega=e^{2\pi i/d}$, one has
$\operatorname{Im}\tau_b<0$. In particular $\tau_b$ is neither
$0$ nor any rational real number.
\end{lemma}

\begin{proof}
The classical Rogers--Szeg\H{o} recurrence and values at $-1$ are
\[
H_{n+1}=(1+z)H_n-z(1-q^n)H_{n-1},\qquad
H_{2j}(-1;q)=\poch{q}{q^2}{j},\quad H_{2j+1}(-1;q)=0.
\]
Differentiate the recurrence with respect to $z$ at $z=-1$:
\[
g_{2r+1}(q)=\poch{q}{q^2}{r}
              +(1-q^{2r})g_{2r-1}(q).
\]
Divide by $\poch{q^2}{q^2}{r}$ and induct from $g_1=1$.
This proves the first identity in (5) wherever the denominators
are nonzero. They are all nonzero at $\omega$, since $2r<d$.

Put $E=\poch{\omega^2}{\omega^2}{h}$. Cancellation of the odd
factors in (4) gives $W_b=\poch{\omega^2}{\omega^2}{r}/E$.
Also $H_d(z;\omega)=1+z^d$, so the same derivative identity at
$r=h$ gives $S_h=d/E$. This proves the ratio in (5).

For the principal embedding, write $\theta=\pi/d$. Every factor
has the exact phase giving
\[
t_k=a_k e^{-ik\theta},\qquad
a_k=\prod_{j=1}^k
\frac{\sin((2j-1)\theta)}{\sin(2j\theta)}>0.
\tag{6}
\]
The arguments lie in $[-h\theta,0]\subset(-\pi/2,0]$.
Hence $\operatorname{Re}S_r>0$, proving $g_b\ne0$.
Moreover, since $r<h$,
\[
\operatorname{Im}\big((S_h-S_r)\overline{S_r}\big)
=-\sum_{j=r+1}^{h}\sum_{k=0}^{r}
a_ja_k\sin((j-k)\theta)<0.
\]
Divide by $|S_r|^2$ to obtain $\operatorname{Im}\tau_b<0$.
The nonvanishing conclusions are algebraic, so they transfer to
every primitive embedding. Nonreality transfers as well because
cyclotomic automorphisms commute with complex conjugation.
\end{proof}

There is an additional exact symmetry. For the complementary odd
residual $b'=d-b-1$,
\[
\tau_{b'}=\overline{\frac{\tau_b}{\tau_b-1}},\qquad
|\tau_b-1|^2|\tau_{b'}-1|^2=1.
\tag{7}
\]
Indeed $t_{h-k}=t_h\overline{t_k}$, $|t_h|=1$, and
$S_h-S_r=t_h\overline{S_{h-r-1}}$. The principal-embedding
quantities $|S_j|$ increase strictly, since every pair of terms in
(6) makes an angle smaller than $\pi/2$. Therefore
$|\tau_b-1|\leq1$ precisely when $2b\geq d-1$. This gave an initial
extension to the upper half of the residual range. The arithmetic
argument below removes that restriction entirely.

\section{The new arithmetic step: a fixed half-integer valuation}

Normalize every extension of the $2$-adic valuation by $v_2(2)=1$.
Define integral monic parity polynomials
\[
\He_{2m}(x)=U_m(x^2),\qquad
\He_{2m+1}(x)=xV_m(x^2).
\]
For either family $R=U$ or $R=V$, define
\[
\mathcal W_m^R(y)
=R_m'(y)R_{m-1}(y)-R_m(y)R_{m-1}'(y),
\qquad m\geq1.
\]
It is monic of degree $2m-2$; $\mathcal W_1^R=1$.

\begin{theorem}[Exact valuation of parity-Wronskian roots]
\label{thm:valuation}
For $m\geq2$, every root $y$ of either $\mathcal W_m^U$ or
$\mathcal W_m^V$, in an algebraic closure of $\Q_2$, satisfies
\[
\boxed{v_2(y-1)=\frac12.}
\tag{8}
\]
Equivalently, the lower Newton polygon of
$\mathcal W_m^R(1+z)$ is the single segment from $(0,m-1)$ to
$(2m-2,0)$.
\end{theorem}

\begin{proof}
Two steps of the Hermite recurrence give
\begin{align*}
U_{m+1}&=(y-4m-1)U_m-2m(2m-1)U_{m-1},\\
V_{m+1}&=(y-4m-3)V_m-2m(2m+1)V_{m-1},
\end{align*}
with $U_0=V_0=1$, $U_1=y-1$, and $V_1=y-3$.
Choose $\pi$ with $\pi^2=2$ and put
\[
P_m(X)=\pi^{-m}U_m(1+\pi X),\qquad
Q_m(X)=\pi^{-m}V_m(1+\pi X).
\]
The exact scaled recurrences are
\begin{align*}
P_{m+1}&=(X-2\pi m)P_m-m(2m-1)P_{m-1},\\
Q_{m+1}&=(X-\pi(2m+1))Q_m-m(2m+1)Q_{m-1}.
\end{align*}
Their initial values are $P_0=Q_0=1$, $P_1=X$, $Q_1=X-\pi$.
Thus $P_m,Q_m$ belong to $\Z[\pi][X]$. Modulo the uniformizer
$\pi$ of $\Q_2(\pi)$, both families reduce to
\[
\rho_m(X)=X^{m\bmod2}(X^2+1)^{\lfloor m/2\rfloor}
\quad\text{in }\F_2[X].
\tag{9}
\]
This follows directly from
$\rho_{m+1}=X\rho_m+(m\bmod2)\rho_{m-1}$.

The chain rule yields
\[
\pi^{-2m+2}\mathcal W_m^R(1+\pi X)
=\widetilde R_m'\widetilde R_{m-1}
 -\widetilde R_m\widetilde R_{m-1}',
\tag{10}
\]
where $\widetilde R=P$ or $Q$. Differentiation commutes with
reduction. In characteristic two, the derivative of the even-index
$\rho$ is zero and that of the odd-index $\rho$ is the remaining
power of $X^2+1$. Hence
\[
\pi^{-2m+2}\mathcal W_m^R(1+\pi X)
\equiv (X^2+1)^{m-1}\pmod\pi.
\tag{11}
\]
The polynomial on the left is monic and integral, with a unit constant
coefficient. Every root $X$ has valuation zero: a negative valuation
would make the leading monomial the unique term of least valuation;
a positive valuation would make the constant term the unique such
term. Either is incompatible with a zero sum.
Since $y-1=\pi X$, (8) follows.

For the Newton-polygon statement, write
$\mathcal W_m^R(1+z)=\sum_j c_jz^j$.
Integrality in (10) says
$v_2(c_j)\geq m-1-j/2$; (11) makes the constant coefficient
attain equality, while the leading coefficient is $1$.
These facts give exactly the stated segment.
\end{proof}

\begin{theorem}[Collision exclusion over unramified fields]
\label{thm:unramified}
Let $K\subseteq\Q_2^{\mathrm{nr}}$, the maximal unramified algebraic
extension of $\Q_2$, and let $\tau\in K$. For every $A\geq1$,
\[
\begin{array}{ll}
A\text{ odd}:&
J_{A+1}(x;\tau)\text{ is squarefree}\ \Longleftrightarrow\ \tau\ne0,\\[2pt]
A\text{ even}:&
J_{A+1}(x;\tau)\text{ is squarefree}\ \Longleftrightarrow\ \tau\ne-1/A.
\end{array}
\tag{12}
\]
\end{theorem}

\begin{proof}
Suppose first that $x\ne0$ is a repeated root. Consecutive Hermite
polynomials are coprime by their recurrence, so
$\He_{A-1}(x)\ne0$. Differentiate $J$ and eliminate $\He_A$:
\[
J'=(1+A\tau)\He_A+A(1-\tau)x\He_{A-1},\qquad
(\tau-1)x^2=\tau(A\tau+1).
\tag{13}
\]
If $\tau=1$, the last equality contradicts $A+1\ne0$.
Otherwise
\[
y=x^2=\frac{\tau(A\tau+1)}{\tau-1}\in K.
\tag{14}
\]
For $A=2m-1$ or $A=2m$, respectively,
\begin{align*}
J_{2m}(x;\tau)&=U_m(y)-A(\tau-1)U_{m-1}(y),\\
J_{2m+1}(x;\tau)&=x\bigl(V_m(y)-A(\tau-1)V_{m-1}(y)\bigr).
\end{align*}
Since $x\ne0$, repetition gives a common root of the relevant
polynomial in $y$ and its derivative. Eliminating $A(\tau-1)$
gives $\mathcal W_m^R(y)=0$. For $m=1$ this is impossible because
the Wronskian is $1$. For $m\geq2$, Theorem~\ref{thm:valuation}
contradicts $y\in K$: an unramified extension has integer
valuation group.

It remains to consider zero. If $A$ is odd,
$J(0)=-A\tau\He_{A-1}(0)$, and the even Hermite value is nonzero.
Thus zero is a root only for $\tau=0$; then $J=x\He_A$ has a
double zero. If $A$ is even, $J$ is odd and
$J'(0)=(1+A\tau)\He_A(0)$, giving exactly $\tau=-1/A$.
All nonzero roots have already been treated. The small cases are
$J_2=x^2-\tau$ and $J_3=x(x^2-1-2\tau)$, consistent with (12).
\end{proof}

\begin{corollary}[All odd-root-order overlap guards]\label{cor:allJ}
For every odd $d\geq3$, every primitive $d$th root $\omega$, every
odd $1\leq b<d$, and every $A\geq1$,
\[
g_b\ne0,\qquad \Disc_xJ_{A+1}(x;\tau_b)\ne0.
\tag{15}
\]
\end{corollary}

\begin{proof}
An odd cyclotomic field embeds into an unramified extension of
$\Q_2$. Explicitly, choose $f$ with $d\mid2^f-1$; a residue-field
element of order $d$ has a unique Hensel lift among the roots of
$X^d-1$ in the degree-$f$ unramified extension, since $d$ is odd.
Lemma~\ref{lem:partial} excludes both rational exceptional parameters
in (12). Apply Theorem~\ref{thm:unramified}.
\end{proof}

\section{Removing the static conditions at prime orders}

\begin{theorem}[Prime-power static nonvanishing]\label{thm:static}
Let $d=p^s$ with $p$ an odd prime, and let $0\leq b<p$.
For a primitive $d$th root $\omega$, $H_b(z;\omega)$ is squarefree.
If $b$ is even, it is relatively prime to $1+z^d$.
If $b$ is odd and $H_b=(z+1)U_b$, then $U_b$ is squarefree,
$U_b(-1)\ne0$, and $\Res(1+z^d,U_b)\ne0$.
\end{theorem}

\begin{proof}
The recurrence gives an integral Hermite initial form. Substitute
$q=1+t^2$, $z=-1+tx$ and induct:
\[
H_b(-1+tx;1+t^2)=t^b K_b(x,t),\qquad
K_b\in\Z[x,t],\quad K_b(x,0)=\He_b(x).
\tag{16}
\]
The polynomial $K_b$ is monic in $x$. Taking discriminants gives
\[
D_b(1+t^2)=t^{b(b-1)}\Disc_x K_b(x,t).
\]
Therefore
\[
D_b(q)=(q-1)^{b(b-1)/2}V_b(q),\qquad
V_b\in\Z[q],\qquad V_b(1)=\prod_{j=1}^b j^j.
\tag{17}
\]
Here $\Disc(\He_b)=\prod_{j=1}^b j^j$ is the classical Hermite
discriminant, also obtained from $\He_b'=b\He_{b-1}$ and the
recurrence for consecutive resultants.

If $D_b(\omega)=0$, then $V_b(\omega)=0$, because $\omega\ne1$.
The monic minimal polynomial $\Phi_{p^s}$ divides $V_b$ in
$\Z[q]$. At $q=1$ this forces
$p=\Phi_{p^s}(1)\mid\prod_{j=1}^b j^j$, impossible since $b<p$.
This proves squarefreeness.

All roots of $1+z^d$ are $-\omega^j$, $0\leq j<d$.
For $b=2r+1$, the identity $H_b(-1;q)=0$ gives
$H_b(z;q)=(z+1)U_b(z;q)$ in $\Z[z,q]$.
Dividing (16) by $tx$ gives
\[
U_b(-1+tx;1+t^2)=t^{2r}L_b(x,t),\qquad
L_b(x,0)=\He_b(x)/x.
\]
For each $j\geq0$, substitute
$x_j(t)=(1-(1+t^2)^j)/t\in t\Z[t]$.
It follows, as an identity of integral polynomials, that
\[
U_b(-q^j;q)=(q-1)^rB_{b,j}(q),\qquad
B_{b,j}(1)=(-1)^r b!!.
\tag{18}
\]
The integrality of the quotient follows from division by the monic
polynomial $(q-1)^r$. If $U_b(-\omega^j;\omega)=0$, then
$\Phi_{p^s}\mid B_{b,j}$, forcing $p\mid b!!$, again impossible.
This includes $j=0$ and proves every odd residual condition.

For $b=2r$, the same substitution directly in (16) gives
\[
H_b(-q^j;q)=(q-1)^rC_{b,j}(q),\qquad
C_{b,j}(1)=(-1)^r(2r-1)!!.
\tag{19}
\]
Its value cannot be divisible by $p>b$, proving separation.
For $b=0$ the polynomial is $1$; for $b=1$ the quotient $U_b$ is
$1$. Both boundary cases agree with the argument.
\end{proof}

\section{Composition with the existing local theorem}\label{sec:leading}

We now state the established local expansion in the notation needed
to compose the new results. Equations (20)--(25) are used under
conditions (3), now proved automatically for $d=p^s$ and $b<p$.
Fix odd $d$, $n=Ad+b$, $A\geq1$,
and use
\[
q(\eps)=\omega(1+\eps^2)^{1/d},\qquad q(\eps)^d-1=\eps^2,
\quad F=1+z^d,\quad G=H_b(z;\omega),\quad
\mathcal P_A=\prod_{j=1}^A j^j.
\]
At a separated fixed root $\alpha$ of $F$, the existing theorem gives
\[
H_n(\alpha+\eps x;q(\eps))
=\eps^A(-1)^A d^{A/2}G(\alpha)
 \He_A\!\left(\frac{\sqrt d\,x}{\alpha}\right)
+O(\eps^{A+1}).
\tag{20}
\]
At the simple overlap $-1$ for odd $b$, it gives
\[
H_n(-1+\eps x;q(\eps))
=\eps^{A+1}d^{(A-1)/2}g
 J_{A+1}(\sqrt d\,x;\tau_b)+O(\eps^{A+2}).
\tag{21}
\]
These are the prior version-3 results, not claimed as new here.
The conditions in (3) ensure that the initial roots are distinct
and that the root clusters account for the full degree.

For even $b$, define $M_{\rm sep}=dA(A-1)/2$. The previous
root-cluster calculation gives
\[
D_n(q(\eps))=\eps^{2M_{\rm sep}}
 \bigl(C_{\rm sep}+O(\eps^2)\bigr),
\]
\[
\boxed{
C_{\rm sep}=(-1)^{A(d-1)/2}
d^{dA(A+1)/2}\mathcal P_A^d
\Disc(G)\Res(F,G)^{2A}.}
\tag{22}
\]
For odd $b$, write $G=(z+1)U$, $g=U(-1)$ and
$M_{\rm ov}=dA(A-1)/2+A$. Then
\[
D_n(q(\eps))=\eps^{2M_{\rm ov}}
 \bigl(C_{\rm ov}+O(\eps^2)\bigr),
\]
\begin{align}
\boxed{\begin{aligned}
C_{\rm ov}={}&(-1)^{A(d-1)/2}
d^{dA(A+1)/2+A}\mathcal P_A^{d-1}
\Disc(J_{A+1})\\[-2pt]
&{}\cdot\Disc(U)\Res(F,U)^{2A}g^2.
\end{aligned}}
\tag{23}
\end{align}

For $d=p^s$ and $b<p$, Theorem~\ref{thm:static} makes all static
factors nonzero; Corollary~\ref{cor:allJ} supplies the remaining
odd-case factor. Hence (22) or (23) is nonzero, proving (2).
The prime case gives (1). For $A=0$, Theorem~\ref{thm:static}
directly gives $D_b(\omega)\ne0$.

The parameter conversion is material. If $M$ is the relevant order
and $C$ the coefficient in (22) or (23), then
\[
\lim_{q\to\omega}\frac{D_n(q)}{(q^d-1)^M}=C,
\qquad
\lim_{q\to\omega}\frac{D_n(q)}{(q-\omega)^M}
=C(d\omega^{-1})^M.
\tag{24}
\]
For the polynomial factor $\Phi_d(q)$ itself,
\[
\left.\frac{D_n(q)}{\Phi_d(q)^M}\right|_{q=\omega}
=C\left(\frac{d\omega^{-1}}{\Phi_d'(\omega)}\right)^M\ne0.
\tag{25}
\]
No coefficient in one local parameter has been substituted for a
coefficient in another.
For $d=p^s$, the identity
$(q^d-1)/\Phi_d(q)=q^{d/p}-1$ simplifies (25) to
\[
\boxed{\left.\frac{D_n(q)}{\Phi_d(q)^M}\right|_{q=\omega}
=C(\omega^{d/p}-1)^M.}
\]
At prime order $d=p$, the conversion factor is simply $(\omega-1)^M$.

\section{Examples and immediate consequences}

Equation (1) applies to all odd prime factors, without an exceptional
residual list. Several examples beyond the old $b=1,3$ result are:
\begin{center}
\begin{tabular}{rrrrl}
\toprule
$p$ & $A$ & $b$ & $n$ & Exact $\Phi_p$ order\\
\midrule
7 & 1 & 5 & 12 & $1$\\
7 & 2 & 5 & 19 & $9$\\
11 & 1 & 5 & 16 & $1$\\
11 & 2 & 7 & 29 & $13$\\
13 & 3 & 5 & 44 & $42$\\
11 & 2 & 6 & 28 & $11$\\
\bottomrule
\end{tabular}
\end{center}
These are substitutions into the proven formula, not six independent
claims inferred from experiments. In particular, for every $A\geq0$,
\[
\ord_{\Phi_{11}}D_{11A+7}
=\frac{11A^2-9A}{2}.
\]

There is also a useful formal decomposition. Put
$e_p(n)=p\binom{\lfloor n/p\rfloor}{2}
+\lfloor n/p\rfloor((n\bmod p)\bmod2)$. Then
\[
D_n(q)=
\left(\prod_{\substack{p\leq n\\p\text{ odd prime}}}
\Phi_p(q)^{e_p(n)}\right)R_n(q),\qquad R_n\in\Z[q],
\]
and no odd-prime cyclotomic polynomial divides $R_n$.
This does not identify the remaining composite-order factors or the
noncyclotomic part of $R_n$.

\section{Verification, source comparison, and scope}

The general statements above rest on the proofs. Finite calculations
were used to identify the arithmetic mechanism and to check
normalizations. They are supplied as short, standard-library Python
programs, with exact integer, rational, or finite-field arithmetic.
The partial-sum controls test the cyclotomic identities independently
of the complex sector proof. The static controls reconstruct small
ordinary discriminants and the integral node quotients in
(17)--(19), and include the excluded boundary $b=p$.
The parity-Wronskian argument has been checked by a second derivation:
an integral recurrence proves
$\mathcal W_m^R(1+2u)\equiv2^{m-1}\pmod{2^m}$, which independently
excludes unramified roots. The stronger proof in
Theorem~\ref{thm:valuation} determines their exact valuation.

One independent control reconstructs the actual $H_{12}$ and its first
$q$-derivative in $\Q(\omega)$, where $\omega$ has order $7$.
Using the reciprocal factorization
$H_{12}(z;q)=z^6R(z+z^{-1};q)$ and
$D_{12}=R(2)R(-2)\Disc(R)^2$, it computes the entire local
coefficient and confirms exactly
\[
D_{12}'(\omega)=7\omega^{-1}C_{\rm ov}\ne0
\qquad(A=1,\ b=5).
\]
Both the $q^7-1$ and $\Phi_7(q)$ conversions also pass.
This tests the sign and all powers in a complete case beyond the
previous $b=1,3$ families. The accompanying composition review checks
all boundary cases and the precise prime-power restriction.

The comparison made in this work is bounded and explicit:
\begin{itemize}
\item The classical Gaussian-binomial recurrence, the Hermite
recurrence and discriminant, and the Hermite--Laguerre quadratic
relations are established inputs. They are not originality claims.
\item Sawa and Uchida's published Theorem 3.6 gives discriminants of
$L_m^{(\alpha)}+cL_{m-1}^{(\alpha)}$. Our parity family fits it with
$A+1=2m+s$, $\alpha=s-\tfrac12$ and
$c=A(\tau-1)/(2m)$. Thus the general quasi-Laguerre discriminant
formula is not new. The additional assertion proved here is the
exclusion of its nonzero collision parameters over unramified
$2$-adic fields, and its use for the RS parameters.
\item Grosu and Grosu's Theorem 7 gives Newton-slope bounds for
Wronskian Hermite polynomials at primes
$p>\max\{2,2m-1\}$, with an additional Vandermonde condition.
Their Theorem 8 uses odd moduli. Those stated results do not
directly cover the shifted valuation at $p=2$ in (8).
\item Bonneux, Dunning and Stevens give broad coefficient formulas
and Hermite--Laguerre Wronskian correspondences. The formulas are
relevant prior work; the present short proof does not establish
that no equivalent valuation consequence exists elsewhere.
\item Lubinsky and Saff's 1987 paper is relevant to static
Rogers--Szeg\H{o} zeros and root-of-unity specializations.
The publisher's abstract and the author's publication record were
checked. The complete relevant body was unavailable in this session,
so no assertion of nonoverlap with every theorem in that paper
is made.
\end{itemize}

The exact project increment is therefore secure relative to the
version-3 theorem and the accepted two-residual review: their
unresolved hypotheses have been removed in the ranges stated above.
The stronger arithmetic lemma is a candidate for publication-level
novelty, subject to a complete literature comparison.

For a general odd composite $d$, (15) holds in the full odd residual
range, but the two static conditions remain open here outside
Theorem~\ref{thm:static}. Nothing in this note treats even $d$, gives a
positive-characteristic version of the main theorem, or supplies a
uniform numerical conditioning estimate as $n$ grows.
In particular, the original arbitrary-input FFT target and its
exponent $\kappa$ remain unchanged.

\clearpage
\section*{References and previous research inputs}
\begingroup\small\raggedright
\begin{enumerate}[label={[\arabic*]},leftmargin=2em]
\item \emph{Local ordinary discriminants of Rogers--Szeg\H{o}
polynomials at odd roots of unity}, working note version 3,
11 October 2026, Section 3, equations (1)--(10).
Retained source identity:\par
\nolinkurl{private-source-locator-omitted}.
\item \emph{RS overlap simplicity: independent mathematical review},
accepted two-subfamily theorem, 11 October 2026.
Retained source identity:\par
\nolinkurl{private-source-locator-omitted}.
\item NIST Digital Library of Mathematical Functions,
\href{https://dlmf.nist.gov/18.7}{Section 18.7}
and \href{https://dlmf.nist.gov/18.9}{Section 18.9}.
\item M. Sawa and Y. Uchida,
\emph{Discriminants of classical quasi-orthogonal polynomials
with application to Diophantine equations},
Journal of the Mathematical Society of Japan 71 (2019), 831--860.
\href{https://www.jstage.jst.go.jp/article/jmath/71/3/71_831/_pdf}
{Published primary text}, Theorem 3.6, p.~845.
The preprint is \href{https://arxiv.org/abs/1802.00605}{arXiv:1802.00605}.
\item C. Grosu and C. Grosu,
\emph{The irreducibility of some Wronskian Hermite polynomials},
\href{https://arxiv.org/html/2007.00065v1}{arXiv:2007.00065v1},
Theorems 7--8 and Section 4.
\item N. Bonneux, C. Dunning and M. Stevens,
\emph{Coefficients of Wronskian Hermite polynomials},
\href{https://arxiv.org/abs/1909.03874}{arXiv:1909.03874}.
\item N. Bonneux and M. Stevens,
\emph{Recurrence relations for Wronskian Laguerre polynomials},
\href{https://arxiv.org/abs/1905.12312}{arXiv:1905.12312}.
\item D. S. Lubinsky and E. B. Saff,
\emph{Convergence of Pad\'e approximants of partial theta functions
and the Rogers--Szeg\H{o} polynomials},
Constructive Approximation 3 (1987), 331--361,
\href{https://doi.org/10.1007/BF01890574}{doi:10.1007/BF01890574}.
Publisher abstract and author bibliography inspected; full-body
comparison remains incomplete.
\end{enumerate}
\endgroup
\end{document}
